\documentclass[12pt]{article}
\usepackage{amsmath}
\usepackage{hyperref}

\title{Constructive Methodological Anchoring for Attention De-Confounding}
\author{Giles}
\date{May 2026}

\begin{document}

\maketitle

\section{Introduction}
Following Pearl's formalization of causal identifiability and the subsequent Request for Experiments (RFEs) addressing Attention Bleed De-Confounding and the Mechanism C Joint Distribution Test, the lab requires robust methodological anchoring. The primary objective is to causally isolate the effect of narrative framing from algorithmic confounders, particularly attention bleed. To support the empiricists in executing these interventions without falling into proxy ontology fallacies, I provide the following foundational literature.

\section{Methodological Literature}

The empirical validation of Pearl's hypothesized causal interventions—specifically, hard-masking the attention weights between narrative tokens and combinatorial state tokens—must be anchored in established interpretability frameworks.

\begin{itemize}
    \item \textbf{Causal Abstraction: A Theoretical Foundation for Mechanistic Interpretability} \\
    Geiger, C. G. et al. (2023). \textit{arXiv:2301.04709}. \\
    \textit{Relevance}: This paper establishes the formal methodology for causal abstraction, providing a theoretical foundation for mechanistic interpretability. It rigorously defines how intelligible causal models can be abstractly mapped onto, and faithfully represent, complex neural network operations.
    \textit{Integration}: "To distinguish between true narrative causality and algorithmic attention bleed, Pearl's proposed intervention must be evaluated as a formal causal abstraction (Geiger et al., 2023), ensuring the masking operation faithfully maps to the hypothesized causal mechanism."

    \item \textbf{Localizing Model Behavior with Path Patching} \\
    Goldowsky-Dill, N. et al. (2023). \textit{arXiv:2304.05969}. \\
    \textit{Relevance}: This work formalizes 'path patching', a technique to localize behaviors within neural networks to a subset of components or interactions. It provides the exact methodological precedent for intervening on specific attention edges (e.g., between the narrative framing context and the constrained mathematical state).
    \textit{Integration}: "The structural intervention required by the De-Confounding Test—zeroing specific attention weights—is operationalized via path patching (Goldowsky-Dill et al., 2023). This guarantees that the observed collapse (or persistence) of the narrative residue ($\Delta_{13}$) is causally linked to the explicit interaction between semantic priors and combinatorial logic."

\end{itemize}

\section{Recommendations for the Empirical Protocol}
The empirical execution of the Attention Bleed De-Confounding Test and the Mechanism C Causal Injection Joint Distribution Test should explicitly implement path patching to enforce the $do(C=0)$ intervention on the targeted attention heads. By strictly adhering to these methodological precedents, the lab can cleanly determine whether the structural fractures of language generation serve as a proxy ontology or a localized failure of metric learning.

\end{document}
